\documentclass[11pt,a4paper]{article}
\usepackage{turkos}
\renewcommand{\headlabel}{Roadmap}
\hypersetup{pdftitle={Turk-OS Study Plan: Roadmap}, pdfauthor={Turk-OS Study Plan}}

\begin{document}
\thispagestyle{plain}

\noindent\begin{tikzpicture}
  \path[use as bounding box] (0,0) rectangle (\linewidth,4.2cm);
  \fill[tkNavy, rounded corners=5pt] (0,0) rectangle (\linewidth,4.2cm);
  \fill[tkRed, rounded corners=5pt] (0,0) rectangle (0.2\linewidth,4.2cm);
  \fill[tkRed] (0.15\linewidth,0) rectangle (0.2\linewidth,4.2cm);
  \node[font=\sffamily\bfseries\footnotesize, text=white, anchor=north] at (0.1\linewidth,3.85cm) {PHASES};
  \node[font=\sffamily\bfseries\fontsize{34}{34}\selectfont, text=white] at (0.1\linewidth,2.05cm) {0--15};
  \node[font=\sffamily\bfseries\fontsize{26}{28}\selectfont, text=white, anchor=north west] at (0.23\linewidth,3.8cm) {Turk-OS};
  \node[font=\sffamily\bfseries\Large, text=white, anchor=north west] at (0.23\linewidth,2.65cm) {Roadmap: from zero to a working OS};
  \node[font=\sffamily\small, text=white!78!tkNavy, anchor=south west, text width=0.74\linewidth, align=left] at (0.23\linewidth,0.35cm)
    {A 16-phase study and build plan for a UNIX-like 32-bit x86 kernel, based on five operating-system textbooks};
\end{tikzpicture}

\vspace{0.4em}
\noindent\tkprogress{0}{100}

\tableofcontents

\section{How to use this plan}

You know basic C, you have never written assembly, and you have never built an operating system. That is a normal
starting point: the authors of \LOB were in a similar position when they wrote it. This plan takes you from an empty
folder to \textbf{Turk-OS 1.0}, a small UNIX-like operating system for 32-bit x86 PCs. It boots with GRUB, runs
preemptive multitasking with separate address spaces, loads ELF programs from its own on-disk file system, and drops
you into a shell with pipes and redirection. Everything runs in QEMU first. Real hardware is optional.

The route has sixteen phases, numbered 0 to 15. Each phase has its own PDF in \texttt{docs/pdf/} and ends with a
concrete, testable milestone. The percentages in each phase header (``0\,\% $\rightarrow$ 5\,\%'') estimate how much of
the whole journey is behind you when that milestone is met. Phase 15 ends at 100\,\%, and it also lists optional
tracks for going further.

Every phase document follows the same layout, so you always know where to look. An \emph{at a glance} box gives the
duration, difficulty, prerequisites and result. \emph{Why this phase} explains where it fits. \emph{Concepts} explains
the ideas you must understand before you write code. The \emph{reading plan} lists book chapters with printed page
numbers, each marked \must, \should or \could. \emph{Build it} breaks the work into numbered tasks with starter
snippets and a way to verify each one. The document ends with a debugging table, the milestone checklist, self-check
questions, and stretch goals.

\subsection{The study loop}

The same loop works in every phase. Keep to it, because skipping the reading makes the coding much slower.

\begin{enumerate}
  \item \textbf{Read the theory first}, usually \OSTEP, then the \must items from the other books.
  \item \textbf{Read the practical chapter} in \LOB (when there is one) to see the concrete x86 steps.
  \item \textbf{Look at a real system.} The MINIX lens box points to the matching \OSDI section and source file.
  \item \textbf{Build in small steps.} One task at a time, boot after each change, commit when it works.
  \item \textbf{Verify.} Every task has a \emph{Verify} line. A task is not finished until you have checked it.
  \item \textbf{Write it down.} Keep a lab journal (\texttt{docs/journal.md}) with what you tried, what broke and why.
  \item \textbf{Answer the self-check questions} without looking. If you can't, go back to the reading.
\end{enumerate}

\subsection{Time budget}

The estimates assume about 10--12 hours a week, which adds up to roughly 43 weeks (about ten months). Some phases
will go faster than planned and some much slower; paging (Phase 6) and processes (Phase 8) are where most people get
stuck. What counts is meeting each phase's exit criteria. Never start a phase while the previous milestone is still
shaky, because every later phase builds on top of it.

\section{Your five books and how each one is used}

The five books complement each other. Only \LOB tells you what to type. The others explain why it works and show how
real systems do it. None of them comes with this plan, so get your own copies before you start Phase 0;
Section~\ref{sec:getbooks} says where.

\begin{table}[htbp]
\centering\small\renewcommand{\arraystretch}{1.4}
\begin{tabularx}{\linewidth}{@{}l >{\raggedright\arraybackslash}p{0.2\linewidth} >{\raggedright\arraybackslash}X >{\raggedright\arraybackslash}X@{}}
\toprule
\tkth{Code} & \tkth{Book} & \tkth{Role in this plan} & \tkth{How to read it} \\
\midrule
\LOB & The little book about OS development (78 pp.) & The practical spine. Its 14 short chapters are the build order for Phases 1--14. &
  Read each chapter right before its phase. It uses GRUB Legacy and Bochs; the phase docs give the modern QEMU and GRUB 2 equivalents. \\
\OSTEP & Operating Systems: Three Easy Pieces & Main theory text. Short, friendly chapters built around one problem each. &
  Your first theory read in every phase. Do the homework simulators when a phase suggests it. \\
\OSC & Operating System Concepts, 10th ed. & Encyclopedic reference. Good second explanation and exam-style depth. &
  Use it when an \OSTEP chapter didn't click, or when you want the complete picture of a topic. \\
\MOS & Modern Operating Systems, 4th ed. & Strongest on hardware, I/O, clocks, disks, security and OS design (Ch.~12). &
  Read the hardware and I/O sections closely; they explain the devices you will program. \\
\OSDI & Operating Systems: Design and Implementation, 3rd ed. & The only book with a complete real x86 kernel: MINIX 3, printed line by line. &
  After you understand a topic, read how MINIX implements it. Don't copy it: MINIX is a microkernel, Turk-OS is monolithic. \\
\bottomrule
\end{tabularx}
\caption{Roles of the five books.}
\end{table}

\subsection{Getting your copies}\label{sec:getbooks}

\LOB and \OSTEP are free online from their authors. \OSC, \MOS and \OSDI are commercial textbooks: buy them from the
publisher or a bookshop (a used copy is fine), or borrow them from a library. Table~\ref{tab:editions} lists the
exact editions that the reading plans cite.

\begin{table}[htbp]
\centering\small\renewcommand{\arraystretch}{1.4}
\begin{tabularx}{\linewidth}{@{}l >{\raggedright\arraybackslash}p{0.45\linewidth} >{\raggedright\arraybackslash}X@{}}
\toprule
\tkth{Code} & \tkth{Edition cited in this plan} & \tkth{Where to get it} \\
\midrule
\LOB & Erik Helin, Adam Renberg. \emph{The little book about OS development}, 2015. &
  Free at \url{https://littleosbook.github.io}, as a web page or as a PDF (the \emph{PDF version} link). Page numbers refer to the PDF. \\
\OSTEP & Remzi H.\ Arpaci-Dusseau, Andrea C.\ Arpaci-Dusseau. \emph{Operating Systems: Three Easy Pieces}, version 0.90, Arpaci-Dusseau Books, 2015. &
  Free at \url{https://pages.cs.wisc.edu/~remzi/OSTEP/}, one PDF per chapter, in the current version (1.10). Printed copies are sold on the same page. \\
\OSC & Abraham Silberschatz, Peter Baer Galvin, Greg Gagne. \emph{Operating System Concepts}, 10th ed., Wiley, 2018. ISBN 978-1-119-32091-3. &
  Wiley, a bookshop or a library. \\
\MOS & Andrew S.\ Tanenbaum, Herbert Bos. \emph{Modern Operating Systems}, 4th ed., Pearson, 2015. ISBN 978-0-13-359162-0. &
  Pearson, a bookshop or a library. \\
\OSDI & Andrew S.\ Tanenbaum, Albert S.\ Woodhull. \emph{Operating Systems: Design and Implementation}, 3rd ed., Pearson Prentice Hall, 2006. ISBN 978-0-13-142938-3. &
  A bookshop (often second-hand) or a library. \\
\bottomrule
\end{tabularx}
\caption{The editions this plan cites, and where to get them.}
\label{tab:editions}
\end{table}

\begin{warnbox}[Page numbers belong to these editions]
The page ranges in the reading plans are the \emph{printed} page numbers of the editions in Table~\ref{tab:editions}.
They will not match another edition, and a reflowable e-book has no fixed pages at all, so find each reading by the
chapter and section numbers that every row gives. For \OSTEP this is the normal case: each free chapter PDF starts again at page~1, and the
current version is paginated differently from version 0.90. Chapters 1--43, and every section this plan cites, have
the same numbers in both versions, and the page range still tells you how long a reading is. One title changed:
chapter 23, \emph{The VAX/VMS Virtual Memory System}, is now \emph{Complete Virtual Memory Systems}, with VAX/VMS in \S23.1.
\end{warnbox}

\subsection{References outside the five books}

Books explain ideas, but writing a kernel also needs exact hardware facts. Keep these four free references at hand,
and note that \LOB relies on them too. The \textbf{Intel 64 and IA-32 Architectures Software Developer's Manual,
Volume 3A} (System Programming Guide) is the definitive source for protection, segmentation, paging, interrupts and
the TSS; you will read chapters 2--7 in pieces. The \textbf{OSDev Wiki} (\url{https://wiki.osdev.org}) has
practical pages on nearly every device in this plan, plus lists of common mistakes. The \textbf{Multiboot
Specification 0.6.96} defines the contract between GRUB and your kernel. The \textbf{NASM manual} covers assembler
syntax. The MIT \textbf{xv6} teaching kernel is a good optional companion: it is small and UNIX-like, and \OSTEP
refers to it.

\section{Key technical decisions}

These choices hold for the whole plan. Each one follows from the five books, which keeps your reading and your
code in step.

\begin{table}[htbp]
\centering\small\renewcommand{\arraystretch}{1.4}
\begin{tabularx}{\linewidth}{@{}>{\raggedright\arraybackslash}p{0.17\linewidth} >{\raggedright\arraybackslash}p{0.27\linewidth} X@{}}
\toprule
\tkth{Decision} & \tkth{Choice} & \tkth{Why} \\
\midrule
Target CPU & IA-32 (i386), 32-bit protected mode, BIOS boot & \LOB, \OSDI (MINIX 3) and the Intel examples in \OSC\,\S9.6 are all 32-bit x86. Two-level paging is the simplest real MMU to learn. A 64-bit port is a Phase 15 track. \\
Boot loader & GRUB via Multiboot~1; \texttt{qemu -kernel} for quick runs & Writing a boot loader is a detour. You write one 512-byte boot sector in Phase 1 to learn how booting works, then let GRUB do the job. \\
Languages & C (gnu11) plus NASM (Intel syntax) & Assembly only where C cannot work: the entry point, descriptor table loads, interrupt stubs, context switch and entering user mode. That is under 5\,\% of the code. \\
Toolchain & \texttt{i686-elf} cross GCC and binutils, GNU Make, Git & A cross compiler never assumes a Linux target, which prevents a family of confusing bugs (Phase 0). \\
Emulator, debugger & QEMU and GDB (remote stub); Bochs is optional & QEMU is fast, scriptable, and has a monitor for inspecting registers and page tables. \LOB uses Bochs, whose debugger is still handy for descriptor-table bugs. \\
Kernel structure & Monolithic, modular source tree & Simplest to get right for a first kernel. MINIX's microkernel is studied as a contrast, with an optional experiment in Phase 15. \\
On-disk file system & MINIX V3 layout with 1\,KiB blocks, named TurkFS & The format described in \OSDI\,\S5.6, with source. Linux's \texttt{mkfs.minix -3} and \texttt{fsck.minix} create and check your images, and Linux can even mount them. \\
Testing & In-kernel \texttt{ktest} suites, headless QEMU, \texttt{isa-debug-exit} & \texttt{make test} boots, runs every test and exits with a pass or fail code, so regressions show up at once (Phase 7). \\
\bottomrule
\end{tabularx}
\caption{Technical decisions for Turk-OS.}
\end{table}

\begin{decisionbox}{Why monolithic, when one of your books is about a microkernel?}
In a monolithic kernel, drivers, the file system and the scheduler all run in ring 0 and call each other as plain C
functions. In a microkernel like MINIX 3 they run as separate user-space processes that talk by message passing
(\OSDI\,\S1.5.5, \S2.5.1; \MOS\,\S1.7). Microkernels isolate faults better, but they force you to build IPC,
servers and a protocol for every operation before anything works. For a first kernel, plain function calls keep
the path from idea to running code short. You will still read the MINIX code in every phase, and noticing where the
two designs differ is one of the best ways to understand both.
\end{decisionbox}

\section{The phase map}

The sixteen phases fall into four stages. Arrows show the build order. Each phase depends on everything before it.

\begin{figure}[H]
\centering
\begin{tikzpicture}[
  ph/.style={draw, rounded corners=3pt, line width=0.7pt, minimum width=3.45cm, minimum height=1.3cm,
    text width=3.2cm, align=center, font=\sffamily\scriptsize, inner sep=3pt},
  s1/.style={ph, draw=tkRed!80!black, fill=tkRed!7},
  s2/.style={ph, draw=tkAmber!90!black, fill=tkAmber!9},
  s3/.style={ph, draw=tkBlue!80!black, fill=tkBlue!7},
  s4/.style={ph, draw=tkGreen!80!black, fill=tkGreen!8},
  flow/.style={-{Stealth[length=2mm]}, line width=0.8pt, draw=tkNavy!70},
]
\node[s1] (p0) at (0,0)      {\textbf{P0\enspace Foundations}\\toolchain, C for kernels};
\node[s1] (p1) at (4.1,0)    {\textbf{P1\enspace Assembly \& boot}\\NASM, BIOS, Multiboot};
\node[s1] (p2) at (8.2,0)    {\textbf{P2\enspace Getting to C}\\VGA, serial, kprintf};
\node[s1] (p3) at (12.3,0)   {\textbf{P3\enspace GDT \& IDT}\\CPU exceptions};
\node[s1] (p4) at (12.3,-2)  {\textbf{P4\enspace Hardware IRQs}\\PIC, timer, keyboard};
\node[s2] (p5) at (8.2,-2)   {\textbf{P5\enspace Physical memory}\\page-frame allocator};
\node[s2] (p6) at (4.1,-2)   {\textbf{P6\enspace Paging}\\higher-half kernel};
\node[s2] (p7) at (0,-2)     {\textbf{P7\enspace Kernel heap}\\kmalloc, lists, tests};
\node[s3] (p8) at (0,-4)     {\textbf{P8\enspace Processes}\\context switch, scheduler};
\node[s3] (p9) at (4.1,-4)   {\textbf{P9\enspace Synchronization}\\locks, semaphores};
\node[s3] (p10) at (8.2,-4)  {\textbf{P10\enspace User mode}\\ring 3, system calls};
\node[s3] (p11) at (12.3,-4) {\textbf{P11\enspace Process API}\\fork, exec, wait, COW};
\node[s4] (p12) at (12.3,-6) {\textbf{P12\enspace Storage}\\ATA driver, buffer cache};
\node[s4] (p13) at (8.2,-6)  {\textbf{P13\enspace File systems}\\VFS, initrd, TurkFS};
\node[s4] (p14) at (4.1,-6)  {\textbf{P14\enspace Userland}\\libc, shell, pipes};
\node[s4] (p15) at (0,-6)    {\textbf{P15\enspace Turk-OS 1.0}\\hardening, release};
\foreach \a/\b in {p0/p1,p1/p2,p2/p3,p3/p4,p4/p5,p5/p6,p6/p7,p7/p8,p8/p9,p9/p10,p10/p11,p11/p12,p12/p13,p13/p14,p14/p15}
  \draw[flow] (\a) -- (\b);
\end{tikzpicture}

\vspace{0.6em}
{\sffamily\scriptsize
\begin{tabular}{@{}l@{\hspace{3em}}l@{}}
\tikz\fill[tkRed!60] (0,0) rectangle (0.25,0.25); Stage I: Bare metal (0--28\,\%) &
\tikz\fill[tkAmber!70] (0,0) rectangle (0.25,0.25); Stage II: Memory (28--48\,\%) \\[2pt]
\tikz\fill[tkBlue!60] (0,0) rectangle (0.25,0.25); Stage III: Processes and protection (48--77\,\%) &
\tikz\fill[tkGreen!60] (0,0) rectangle (0.25,0.25); Stage IV: Persistence and userland (77--100\,\%) \\
\end{tabular}}
\caption{The sixteen phases, grouped into four stages.}
\end{figure}

\textbf{Stage I, Bare metal (Phases 0--4).} You learn the tools, enough assembly to be dangerous, and how a PC boots.
By the end the kernel runs C code, prints to the screen and the serial port, survives CPU exceptions with a readable
error report, and reacts to the timer and the keyboard through a small built-in command monitor.

\textbf{Stage II, Memory (Phases 5--7).} The kernel learns which RAM exists and hands it out page by page, turns on
paging, moves itself to the top gigabyte of the address space, and gets a \texttt{kmalloc} heap. You also set up the
automated test harness that protects you from regressions for the rest of the project.

\textbf{Stage III, Processes and protection (Phases 8--11).} Multitasking arrives: kernel threads, context switches,
a preemptive scheduler, then the locks that make concurrency safe. Programs move to ring 3 with their own address
spaces and talk to the kernel through system calls. The stage ends with the UNIX process model: \texttt{fork} with
copy-on-write, \texttt{exec}, \texttt{wait} and \texttt{exit}.

\textbf{Stage IV, Persistence and userland (Phases 12--15).} A disk driver and buffer cache give you storage, a VFS
and TurkFS give you files, and a C library, shell and utilities give you a usable system. The last phase hardens,
measures and documents it, tags version 1.0, and lays out the optional tracks.

\section{Phase summary}

\begin{center}
\small\renewcommand{\arraystretch}{1.35}\setlength{\tabcolsep}{4pt}
\begin{longtable}{@{}c >{\raggedright\arraybackslash}p{0.24\linewidth} >{\raggedright\arraybackslash}p{0.43\linewidth} c c@{}}
\toprule
\tkth{Phase} & \tkth{Title} & \tkth{You finish with} & \tkth{Weeks} & \tkth{Progress} \\
\midrule
\endfirsthead
\toprule
\tkth{Phase} & \tkth{Title} & \tkth{You finish with} & \tkth{Weeks} & \tkth{Progress} \\
\midrule
\endhead
\bottomrule
\endlastfoot
0  & Foundations and toolchain & Cross compiler, QEMU, GDB and a Git repo; C warm-up library with unit tests & 2--3 & 0\,\%\,$\rightarrow$\,5\,\% \\
1  & x86 assembly and booting & A 512-byte boot sector; a Multiboot kernel that GRUB and QEMU boot & 2--3 & 5\,\%\,$\rightarrow$\,10\,\% \\
2  & Getting to C: screen, serial, kprintf & VGA text driver, serial logging, \texttt{kprintf}, \texttt{panic} & 1--2 & 10\,\%\,$\rightarrow$\,16\,\% \\
3  & Segmentation, IDT and exceptions & Your own GDT and IDT; register dump on every CPU exception & 2 & 16\,\%\,$\rightarrow$\,22\,\% \\
4  & Hardware interrupts & PIC, 100\,Hz timer, keyboard driver, kernel command monitor & 2 & 22\,\%\,$\rightarrow$\,28\,\% \\
5  & Physical memory management & Memory map parsing, bitmap page-frame allocator, first kernel tests & 1--2 & 28\,\%\,$\rightarrow$\,34\,\% \\
6  & Paging and the higher-half kernel & Paging on, kernel at \texttt{0xC0000000}, VMM API, page-fault reports & 2--3 & 34\,\%\,$\rightarrow$\,42\,\% \\
7  & Kernel heap and data structures & \texttt{kmalloc}/\texttt{kfree}, intrusive lists, \texttt{make test} in headless QEMU & 1--2 & 42\,\%\,$\rightarrow$\,48\,\% \\
8  & Processes and scheduling & Kernel threads, context switch, preemptive round-robin, sleep & 3 & 48\,\%\,$\rightarrow$\,56\,\% \\
9  & Synchronization & Spinlocks, wait queues, mutexes, semaphores; blocking keyboard input & 2 & 56\,\%\,$\rightarrow$\,62\,\% \\
10 & User mode and system calls & TSS, ring 3, \texttt{int 0x80}, ELF loader, first user programs & 3 & 62\,\%\,$\rightarrow$\,70\,\% \\
11 & Process API & \texttt{fork} (copy-on-write), \texttt{exec}, \texttt{wait}, \texttt{exit}, \texttt{brk}, basic signals & 3 & 70\,\%\,$\rightarrow$\,77\,\% \\
12 & Storage & ATA PIO driver, block layer, MBR partitions, buffer cache & 2 & 77\,\%\,$\rightarrow$\,83\,\% \\
13 & File systems & VFS, file descriptors, tar initrd, writable TurkFS & 3--4 & 83\,\%\,$\rightarrow$\,90\,\% \\
14 & Userland & Turk-libc, pipes, tty line discipline, \texttt{tsh} shell, coreutils & 3 & 90\,\%\,$\rightarrow$\,96\,\% \\
15 & Hardening and release & Fuzzing, permissions, benchmarks, docs, v1.0 ISO; optional tracks & 3--4+ & 96\,\%\,$\rightarrow$\,100\,\% \\
\end{longtable}
\end{center}

\section{Timeline}

The bars in Figure~\ref{fig:timeline} place each phase on a calendar at 10--12 hours a week. Stage colors match the phase map. Expect
Phases 6, 8 and 13 to run over; that is normal, and the rest of the schedule usually absorbs it.

\begin{figure}[htbp]
\centering
\begin{tikzpicture}[x=0.265cm, y=0.37cm]
  \foreach \w in {0,4,...,44} {
    \draw[tkGray!25] (\w,0.7) -- (\w,-15.6);
    \node[tknote, above] at (\w,0.7) {\w};
  }
  \node[tknote, anchor=south west] at (45,0.7) {weeks};
  \foreach \n/\s/\e/\c [count=\i from 0] in {
      P0 Foundations/0/3/tkRed, P1 Assembly and boot/3/6/tkRed, P2 Getting to C/6/8/tkRed,
      P3 GDT and IDT/8/10/tkRed, P4 Hardware IRQs/10/12/tkRed,
      P5 Physical memory/12/14/tkAmber, P6 Paging/14/17/tkAmber, P7 Kernel heap/17/19/tkAmber,
      P8 Processes/19/22/tkBlue, P9 Synchronization/22/24/tkBlue, P10 User mode/24/27/tkBlue,
      P11 Process API/27/30/tkBlue,
      P12 Storage/30/32/tkGreen, P13 File systems/32/36/tkGreen, P14 Userland/36/39/tkGreen,
      P15 Turk-OS 1.0/39/43/tkGreen} {
    \node[font=\sffamily\scriptsize, anchor=east] at (-0.4,-\i) {\n};
    \fill[\c!70, rounded corners=1pt] (\s,-\i-0.32) rectangle (\e,-\i+0.32);
  }
  \draw[tkRed, line width=0.8pt, dashed] (43,0.7) -- (43,-15.6) node[below, font=\sffamily\scriptsize\bfseries, text=tkRed] {v1.0};
\end{tikzpicture}
\caption{Planned timeline at 10--12 hours per week. Treat it as a guide, not a deadline.}
\label{fig:timeline}
\end{figure}

\section{What Turk-OS 1.0 looks like}

Figure~\ref{fig:arch} shows the finished system and the phase that builds each part. It is worth coming back to this
picture whenever you lose track of where a piece fits.

\begin{figure}[htbp]
\centering
\begin{tikzpicture}[
  layer/.style={draw=tkNavy!60, rounded corners=2pt, minimum height=1.05cm, font=\sffamily\scriptsize,
    align=center, inner sep=2pt, line width=0.6pt, anchor=north west},
  ph/.style={text=tkRed, font=\sffamily\scriptsize\bfseries},
]
\node[layer, fill=tkGreen!10, minimum width=4.2cm, text width=4.0cm] at (0,0)    {init, \texttt{tsh} shell, coreutils\\{\color{tkRed}\bfseries P11, P14}};
\node[layer, fill=tkGreen!10, minimum width=4.2cm, text width=4.0cm] at (4.35,0) {Turk-libc: stdio, malloc, string\\{\color{tkRed}\bfseries P14}};
\node[layer, fill=tkGreen!10, minimum width=4.3cm, text width=4.1cm] at (8.7,0)  {your own ELF programs\\{\color{tkRed}\bfseries P10, P11}};
\node[layer, fill=tkRed!8, minimum width=13cm, text width=12.8cm, minimum height=0.8cm] at (0,-1.25)
  {System call interface: \texttt{int 0x80}, argument checks, \texttt{copy\_from\_user} / \texttt{copy\_to\_user}\quad{\color{tkRed}\bfseries P10}};
\node[layer, fill=tkBlue!7, minimum width=3.1cm, text width=2.95cm, minimum height=1.25cm] at (0,-2.25)    {Processes and scheduler\\fork, exec, signals\\{\color{tkRed}\bfseries P8, P11}};
\node[layer, fill=tkBlue!7, minimum width=3.1cm, text width=2.95cm, minimum height=1.25cm] at (3.3,-2.25)  {Memory\\PMM, VMM, heap, COW\\{\color{tkRed}\bfseries P5--P7, P11}};
\node[layer, fill=tkBlue!7, minimum width=3.1cm, text width=2.95cm, minimum height=1.25cm] at (6.6,-2.25)  {VFS\\tarfs, TurkFS, devfs, pipes\\{\color{tkRed}\bfseries P13, P14}};
\node[layer, fill=tkBlue!7, minimum width=3.1cm, text width=2.95cm, minimum height=1.25cm] at (9.9,-2.25)  {Synchronization\\spinlocks, mutexes, sems\\{\color{tkRed}\bfseries P9}};
\node[layer, fill=tkAmber!10, minimum width=13cm, text width=12.8cm, minimum height=0.8cm] at (0,-3.7)
  {Drivers: VGA, serial, keyboard, PIT, tty, ATA, RAM disk, block layer, buffer cache\quad{\color{tkRed}\bfseries P2, P4, P12, P14}};
\node[layer, fill=tkAmber!10, minimum width=13cm, text width=12.8cm, minimum height=0.8cm] at (0,-4.7)
  {i386 layer: Multiboot entry, GDT, IDT and ISR stubs, PIC, TSS, paging, context switch\quad{\color{tkRed}\bfseries P1--P4, P6, P8, P10}};
\node[layer, fill=tkGray!12, minimum width=13cm, text width=12.8cm, minimum height=0.8cm] at (0,-5.9)
  {PC hardware (QEMU): CPU, RAM, VGA, 8259 PIC, 8254 PIT, 8042 keyboard controller, IDE disk, 16550 UART};
\draw[decorate, decoration={brace, amplitude=5pt, mirror}, tkGray, line width=0.7pt] (-0.15,0) -- (-0.15,-1.05)
  node[midway, left=6pt, font=\sffamily\scriptsize\bfseries, text=tkGray, align=right] {user\\space\\(ring 3)};
\draw[decorate, decoration={brace, amplitude=5pt, mirror}, tkGray, line width=0.7pt] (-0.15,-1.25) -- (-0.15,-5.5)
  node[midway, left=6pt, font=\sffamily\scriptsize\bfseries, text=tkGray, align=right] {kernel\\(ring 0)};
\draw[decorate, decoration={brace, amplitude=5pt, mirror}, tkGray, line width=0.7pt] (-0.15,-5.9) -- (-0.15,-6.7)
  node[midway, left=6pt, font=\sffamily\scriptsize\bfseries, text=tkGray, align=right] {hardware};
\end{tikzpicture}
\caption{Layers of Turk-OS 1.0, with the phase that builds each part.}
\label{fig:arch}
\end{figure}

The source tree grows into the layout below. You don't need to create it all on day one; each phase document says
which folders it adds.

\begin{lst}{plain}{Turk-OS/ (final layout)}
Turk-OS/
|-- docs/                 this study plan (LaTeX), journal.md, architecture notes
|-- boot/                  loader.s (Multiboot entry), grub.cfg
|-- kernel/
|   |-- arch/i386/         gdt.c idt.c isr.s irq.c pic.c pit.c tss.c paging.c switch.s
|   |-- mm/                pmm.c vmm.c heap.c
|   |-- proc/              task.c sched.c fork.c exec.c signal.c
|   |-- sync/              spinlock.c mutex.c semaphore.c waitq.c
|   |-- fs/                vfs.c file.c tarfs.c turkfs.c devfs.c pipe.c
|   |-- drivers/           vga.c serial.c keyboard.c tty.c ata.c ramdisk.c blockdev.c bcache.c
|   |-- syscall/           syscall.c uaccess.c
|   `-- kmain.c
|-- lib/                   libk: string.c printf.c bitmap.c ringbuf.c list.h
|-- include/               kernel headers
|-- user/
|   |-- libc/              Turk-libc
|   |-- init/ sh/ bin/     init, tsh, coreutils
|   `-- user.ld
|-- tests/                 ktest suites
|-- tools/                 mkturkfs.c and helper scripts
|-- link.ld
`-- Makefile
\end{lst}

\section{Risks, and what to do when you are stuck}

Everyone who writes a kernel hits days where the machine silently reboots and nothing makes sense. The table lists
the predictable risks. The protocol after it is how you get unstuck without losing a week.

\begin{center}
\small\renewcommand{\arraystretch}{1.35}\setlength{\tabcolsep}{4pt}
\begin{longtable}{@{}>{\raggedright\arraybackslash}p{0.33\linewidth} c >{\raggedright\arraybackslash}p{0.5\linewidth}@{}}
\toprule
\tkth{Risk} & \tkth{Likelihood} & \tkth{Mitigation} \\
\midrule
\endfirsthead
\toprule
\tkth{Risk} & \tkth{Likelihood} & \tkth{Mitigation} \\
\midrule
\endhead
\bottomrule
\endlastfoot
Triple faults (instant reboot) with no output & High & Serial logging from Phase 2 on; run QEMU with \texttt{-d int,cpu\_reset -no-reboot}; GDB; small commits so \texttt{git bisect} can find the change that broke it. \\
Assembly learning curve & High & Phase 1 is dedicated to it. Read compiler output instead of writing assembly from scratch, and keep assembly to tiny stubs. \\
Memory corruption that shows up far from its cause & High & Assertions everywhere, magic numbers in headers, guard pages, \texttt{heap\_check()}, deterministic tests. \\
Stuck on one bug for days & High & Use the protocol below. Time-box the bug, then change approach. \\
Scope creep (GUI or networking too early) & Medium & Exit criteria are the definition of done. New ideas go on a ``later'' list for Phase 15. \\
Toolchain trouble (cross-GCC build, Flatpak sandbox) & Medium & Phase 0 checklist, a host terminal, documented fallback compiler flags. \\
Outdated instructions in \LOB & High & The phase docs give current equivalents (\texttt{grub2-mkrescue}, QEMU); the OSDev Wiki has errata. \\
Losing motivation & Medium & A visible milestone every 2--3 weeks. Record a short demo at each milestone and keep the journal. \\
\end{longtable}
\end{center}

\begin{tipbox}[The getting-unstuck protocol]
\begin{enumerate}
  \item \textbf{Reproduce it reliably}, then make it smaller: comment code out until you find the shortest path to the bug.
  \item \textbf{Look, don't guess.} Use the serial log, \texttt{-d int}, the QEMU monitor (\texttt{info registers}, \texttt{info mem}), and GDB breakpoints.
  \item \textbf{Write the hypothesis down} in your journal before you change code. Change one thing at a time.
  \item \textbf{Re-read the section} of the Intel manual or the book chapter for the exact hardware rule involved.
  \item \textbf{Compare with MINIX} (\OSDI source) or xv6: how does a working kernel do the same step?
  \item \textbf{After two sessions without progress}, ask for help (OSDev forum, a classmate, your advisor) and include your journal notes. Explaining the problem often solves it.
\end{enumerate}
\end{tipbox}

\section{Master milestone tracker}

Print this page and tick each box when that phase's exit criteria are fully met.

\begin{checklist}
  \item \textbf{P0}\enspace \texttt{i686-elf-gcc}, NASM, QEMU and GDB work; repo created; C warm-up library passes its tests.
  \item \textbf{P1}\enspace Boot sector prints a message; Multiboot kernel boots via \texttt{-kernel} and via a GRUB ISO.
  \item \textbf{P2}\enspace Colored banner on screen, same log on the serial port, \texttt{kprintf} and \texttt{panic} tested.
  \item \textbf{P3}\enspace Own GDT and IDT loaded; divide-by-zero shows a register dump; \texttt{int3} returns cleanly.
  \item \textbf{P4}\enspace Timer ticks at 100\,Hz; keyboard input with shift and caps; monitor commands run.
  \item \textbf{P5}\enspace Memory map printed; frame allocator passes its tests from 16\,MiB to 3\,GiB of RAM.
  \item \textbf{P6}\enspace Kernel runs at \texttt{0xC0000000}; null dereference produces a clean page-fault report.
  \item \textbf{P7}\enspace 10\,000-operation heap stress test passes; \texttt{make test} runs headless.
  \item \textbf{P8}\enspace Several kernel threads preempt each other; \texttt{sleep} works; no leaks over 1000 task cycles.
  \item \textbf{P9}\enspace Race demo fixed by locks; bounded buffer and dining philosophers run without deadlock.
  \item \textbf{P10}\enspace Ring-3 program prints through \texttt{write}; faulty programs are killed, the kernel survives.
  \item \textbf{P11}\enspace \texttt{fork}/\texttt{exec}/\texttt{wait} chain works; copy-on-write verified; no leaks over 1000 forks.
  \item \textbf{P12}\enspace QEMU disk identified; data survives reboot; partitions detected; buffer cache in use.
  \item \textbf{P13}\enspace Files written by Turk-OS pass \texttt{fsck.minix} on the host; file descriptor semantics tested.
  \item \textbf{P14}\enspace Boots to the \texttt{tsh} prompt; pipes, redirection and Ctrl+C work; coreutils present.
  \item \textbf{P15}\enspace One hour of fuzzing without a panic; benchmarks and docs written; v1.0 tagged.
\end{checklist}

\booklegend
\end{document}
